\section{Related Work}
\label{sec:relatedwork}
\subsection{Intrusion Detection Systems for In-Vehicle Networks}
Early research on in-vehicle network security focused predominantly on the Controller Area Network (CAN) protocol, which has been the dominant communication standard in automotive systems for decades. A wide range of IDS approaches have been proposed for CAN: rule-based methods that flag violations of message periodicity and identifier frequency~\cite{song2016intrusion}, statistical approaches based on entropy~\cite{muter2011entropy}, and machine learning methods spanning support vector machines~\cite{bari2023intrusion}, random forests~\cite{caivano2024marea}, recurrent neural networks~\cite{taylor2016anomaly}, and deep neural networks~\cite{kang2016intrusion}. Survey studies have systematically categorized these approaches and identified key limitations, including the substantial re-engineering required to port a given detection scheme to a different vehicle make or model, since CAN message implementations and specifications are largely vehicle-specific~\cite{sharmin2024benchmarking, lampe2023survey}.

With the adoption of Automotive Ethernet as the backbone for high-bandwidth in-vehicle communication, attention has begun to shift toward Ethernet-specific security. Several studies have examined vulnerability analysis and attack modeling for Automotive Ethernet~\cite{de2024systematic, fotouhi2023assessing}, and initial IDS proposals have framed intrusion detection as a traffic classification problem, relying on packet-level features extracted directly from payloads~\cite{jeong2021convolutional}, packet byte features homogenized to a fixed length and reduced by feature selection~\cite{shibly2023feature}, or protocol-specific header fields processed by a sequential classifier~\cite{alkhatib2021some}. However, the majority of existing Automotive Ethernet IDS proposals rely on supervised learning, requiring labeled examples of each attack type during training~\cite{jeong2021convolutional,alkhatib2021some}. This dependency is problematic in practice, as the space of possible attacks against Automotive Ethernet is large and continuously evolving, making comprehensive labeled dataset collection infeasible. In contrast, our method requires no labeled attack data, learning exclusively from normal traffic and detecting anomalies as deviations from the learned normal behavior.

Table~\ref{tab:relatedwork_ae} summarizes these Automotive Ethernet IDS approaches alongside the proposed method. The proposed method is the only one of the compared approaches that is simultaneously label-free in training and protocol-agnostic: AERO removes the labeling requirement but still depends on protocol-informed feature engineering, while the supervised and semi-supervised approaches depend on both labels and protocol-specific representations. Labels re-enter only when the detection threshold is selected, which applies to the proposed method as well (Section~\ref{sec:limitations}).

\begin{table*}[!tbp]
\centering
\caption{Qualitative comparison of existing Automotive Ethernet IDS approaches. Label Req.: fraction of labeled attack data required for training.}
\label{tab:relatedwork_ae}
\footnotesize
\begin{tabular}{lcccc}
\toprule
\textbf{Method} & \textbf{Paradigm} & \textbf{Label Req.} & \textbf{Feature} & \textbf{Protocol-Agnostic} \\
\midrule
Jeong et al.~\cite{jeong2021convolutional} & Supervised & 100\% & Protocol-specific (AVTP payload) & No \\
Alkhatib et al.~\cite{alkhatib2021some} & Supervised & 100\% & Protocol-specific (SOME/IP header fields) & No \\
Shibly et al.~\cite{shibly2023feature} & Semi-supervised & $\sim$10\% & Padded packet bytes + feature selection & No \\
AERO~\cite{jeong2023aero} & Unsupervised & 0\% & Engineered, protocol-informed (multimodal) & No \\
\textbf{Proposed} & \textbf{Self-supervised} & \textbf{0\%} & \textbf{Raw byte} & \textbf{Yes} \\
\bottomrule
\end{tabular}
\end{table*}

\subsection{Anomaly Detection in Network Traffic}
Anomaly-based network intrusion detection has a long research history in the broader network security community. Classical approaches based on statistical modeling, such as principal component analysis (PCA) and one-class support vector machines (OC-SVM), establish a model of normal traffic and flag deviations as anomalies~\cite{chandola2009anomaly}. These methods are well-suited to unsupervised settings but suffer from limited representational capacity when applied to high-dimensional, variable-length packet data.

Deep learning has substantially advanced the state of the art in network anomaly detection. Autoencoder-based methods train a neural network to reconstruct normal traffic and use the reconstruction error as an anomaly score~\cite{mirsky2018kitsune}. Variants including variational autoencoders (VAE)~\cite{kingma2014auto} and adversarially trained models~\cite{schlegl2017anogan} have been proposed to improve the quality of the learned normal distribution. Recurrent architectures such as long short-term memory (LSTM) networks have been applied to model temporal dependencies in traffic sequences~\cite{kim2016lstm}, while convolutional approaches have demonstrated effectiveness at capturing local patterns in packet payload bytes, albeit after protocol-specific extraction~\cite{jeong2021convolutional}. More recently, Transformer-based models have been explored for network traffic analysis, leveraging self-attention to capture long-range dependencies across packet sequences~\cite{lin2022bert}.

Despite these advances, existing network anomaly detection methods are predominantly evaluated on enterprise or IoT network datasets such as KDD Cup 99, NSL-KDD, and CICIDS~\cite{tavallaee2009detailed, sharafaldin2018toward}, which differ substantially from Automotive Ethernet traffic in terms of protocol composition, traffic regularity, and attack characteristics. Our work addresses this gap by targeting Automotive Ethernet specifically and by designing an architecture that accounts for the unique properties of this domain.

\subsection{Self-Supervised Learning for Anomaly Detection}
Self-supervised learning has emerged as a powerful paradigm for representation learning from unlabeled data. Early approaches defined pretext tasks such as rotation prediction~\cite{gidaris2018unsupervised} and contrastive instance discrimination~\cite{chen2020simclr} to learn transferable representations without manual annotation. Masked autoencoding, introduced in the vision domain by the Masked Autoencoder (MAE)~\cite{he2022masked} and in the language domain by BERT~\cite{devlin2019bert}, trains a model to reconstruct randomly masked portions of its input, forcing it to internalize the generative structure of the data. This approach has subsequently been extended to time-series~\cite{cheng2026timemae}, graph-structured data~\cite{hou2022graphmae}, and multimodal inputs~\cite{bachmann2022multimae}.

In the context of anomaly detection, self-supervised methods offer a natural fit: a model trained to reconstruct or predict normal data will produce elevated reconstruction errors when presented with anomalous inputs that deviate from learned patterns. This principle has been applied to anomaly detection in industrial sensor data~\cite{audibert2020usad}, medical time series~\cite{zhang2022maefe}, and video surveillance~\cite{liu2018future}. For network intrusion detection, several recent works have explored self-supervised pretraining as a means of reducing labeled data requirements~\cite{lin2022bert}, though applications to raw packet byte sequences in automotive settings remain limited.

Our work extends masked autoencoding to the Automotive Ethernet domain, introducing a dual-stream architecture that distinguishes between packet content features and inter-packet byte-difference features. Unlike prior works that apply a single homogeneous encoder to network traffic, our design explicitly accounts for the complementary nature of global inter-packet sequential structure and per-byte temporal dynamics derived from inter-packet differences, demonstrating that this separation leads to measurably superior anomaly detection performance.

\subsection{Transformer and CNN Hybrid Architectures}
The complementary strengths of Transformer and CNN architectures have motivated a growing body of research on hybrid designs. Transformers excel at modeling long-range dependencies through the self-attention mechanism~\cite{vaswani2017attention}, while CNNs provide strong inductive biases for capturing local patterns through shared convolutional kernels~\cite{lecun1998gradient}. Hybrid architectures that combine both have demonstrated superior performance over purely Transformer or CNN-based designs in domains including image recognition~\cite{wu2021cvt}, speech processing~\cite{gulati2020conformer}, and time-series forecasting~\cite{zhang2023novel, el2025hybrid}.

Hybrid Transformer-CNN models have been proposed in the network traffic analysis domain for encrypted traffic classification~\cite{shi2023bfcn} and flow-level anomaly detection~\cite{wang2025self}. However, to our knowledge, no prior work has applied a hybrid Transformer-CNN architecture to raw packet byte sequence modeling for Automotive Ethernet intrusion detection. Our ablation studies empirically validate the design choice of combining a Transformer for global inter-packet content modeling with a residual CNN for per-byte temporal dynamics from inter-packet byte differences, demonstrating consistent performance advantages over single-encoder alternatives across diverse attack scenarios.